\input{../tex-inputs/latex-korrekturansicht-vorspann}

\section[Arthur Schnitzler an Gustav Schwarzkopf, 4. 1. 1925]{L04353 Arthur Schnitzler an Gustav Schwarzkopf, 4. 1. 1925}
\nopagebreak\mylabel{L04353v}
\rehead{ }\normalsize\beginnumbering\briefempfaengerindex{Schwarzkopf, Gustav@\textsc{Schwarzkopf, Gustav}!zzzSchnitzler, Arthur@\emph{von Arthur Schnitzler}!1925-01-041@{4. 1. 1925}|(be}
\toendnotes[C]{\smallbreak\pagebreak[2]}
\correspDesc{Versand  durch Arthur Schnitzler am 4. 1. 1925 in Wien
\newline{}Übermittlung  am 5. 1. 1925 in Wien
\newline{}Erhalt  durch Gustav Schwarzkopf im Zeitraum [4. 1. 1925 – 7. 1. 1925?] in Wien}\toendnotes[C]{\smallbreak}
\Standort{DLA, A:Schnitzler, HS.1985.1.1897.}
\physDesc{Postkarte, 459 Zeichen
\newline{}Handschrift: Bleistift, lateinische Kurrent
\newline{}Versand: Stempel: »\nobreak{}\oindex{XVIII., Währing@\textbf{XVIII., Währing}, \emph{Verwaltungsgebiet}|pwk}\textcolor{gray}{18/1} Wien
                                       110, 5. I. 25, 9\nobreak{}«.  }\toendnotes[C]{\smallbreak}\pstart{}{\pb}\label{T_L04353-1v}\edtext{\textcolor{gray}{\textbf{A. S.}}}{\lemma{\textnormal{\emph{A. S.}}}\Cendnote{\textnormal{ovaler Absenderkleber}}}\label{T_L04353-1}\pend{}\pstart{}\textcolor{pink}{\textcolor{gray}{\textbf{WIEN, XVIII.}}}\oindex{XVIII., Währing@\textbf{XVIII., Währing}, \emph{Verwaltungsgebiet}|pw}{}\ledrightnote{\textcolor{pink}{XVIII., Währing}}\pend{}\pstart{}\textcolor{pink}{\textcolor{gray}{\textbf{STERNWARTESTR. 71}}}\oindex{Wien@\textbf{Wien}!XVIII., Währing@\textbf{XVIII., Währing}!Sternwartestraße 71@\textbf{Sternwartestraße 71}, \emph{Wohngebäude}|pw}{}\ledrightnote{\textcolor{pink}{Sternwartestraße 71}}\pend{}{\bigskip}\pstart{}Hr Guſtav Schwarzkopf\pend{}\pstart{}\textcolor{pink}{Wien I}\oindex{I., Innere Stadt@\textbf{I., Innere Stadt}, \emph{Verwaltungsgebiet}|pw}{}\ledrightnote{\textcolor{pink}{I., Innere Stadt}}\pend{}\pstart{}\textcolor{pink}{Tiefer Graben 17}\oindex{Wien@\textbf{Wien}!I., Innere Stadt@\textbf{I., Innere Stadt}!Tiefer Graben 17@\textbf{Tiefer Graben 17}, \emph{Wohngebäude}|pw}{}\ledrightnote{\textcolor{pink}{Tiefer Graben 17}}\pend{}{\bigskip}\vspace{1em}
\pstart
           \raggedleft{}{\pb}\textcolor{pink}{Wien}\oindex{Wien@\textbf{Wien}, \emph{Verwaltungsgebiet}|pw}{}\ledrightnote{\textcolor{pink}{Wien}}{ }4. 1. 925\pend
           \vspace{0.5em}
\pstart
           lieber Guſtav; – Ihr Tel. ist offenbar verdorben ich rief seit
               gestern ein halbes Dutzend Mal an, es kracht unentwegt{[}.{]} Also:
                  \label{K_L04353-1v}\edtext{es k\textcolor{gray}{ä}men in
                  Betracht}{\lemma{\textnormal{\emph{es kämen in
                  Betracht}}}\Cendnote{\textnormal{Es dürfte sich um die Suche
                  für eine psychiatrische Klinik für \textcolor{blue}{Schwarzkopfs}\pwindex{Schwarzkopf, Gustav 7.\,11.\,1853 Wien – 13.\,11.\,1939 ebd.@\textsc{Schwarzkopf, Gustav} (7.\,11.\,1853 Wien – 13.\,11.\,1939 ebd.), \emph{Schriftsteller}|pwk} Bruder \textcolor{blue}{Emil}\pwindex{Schwarzkopf, Emil 17.\,9.\,1851 Wien – 28.\,1.\,1928 ebd.@\textsc{Schwarzkopf, Emil} (17.\,9.\,1851 Wien – 28.\,1.\,1928 ebd.), \emph{Übersetzer, Komponist, Musiklehrer}|pwk} handeln,
                     vgl. A. S.: \emph{Tagebuch}, 18. 11. 1924: »Bei
                        \textcolor{blue}{Gustav}\pwindex{Schwarzkopf, Gustav 7.\,11.\,1853 Wien – 13.\,11.\,1939 ebd.@\textsc{Schwarzkopf, Gustav} (7.\,11.\,1853 Wien – 13.\,11.\,1939 ebd.), \emph{Schriftsteller}|pw}. Sein Bruder \textcolor{blue}{Emil}\pwindex{Schwarzkopf, Emil 17.\,9.\,1851 Wien – 28.\,1.\,1928 ebd.@\textsc{Schwarzkopf, Emil} (17.\,9.\,1851 Wien – 28.\,1.\,1928 ebd.), \emph{Übersetzer, Komponist, Musiklehrer}|pw} bettlägerig; sein schrullenhaft
                     hysterisch-hypochondrisches Wesen ins senil groteske
                  gesteigert.«}}}\label{K_L04353-1}\textcolor{pink}{Rosenhügel}\oindex{Wien@\textbf{Wien}!XIII., Hietzing@\textbf{XIII., Hietzing}!Heilanstalt »Rosenhügel«@\textbf{Heilanstalt »Rosenhügel«}, \emph{Sanatorium}|pw}{}\ledrightnote{\textcolor{pink}{Heilanstalt »Rosenhügel«}}, wo Dr. \textcolor{blue}{\uline{Sölder}}\pwindex{Sölder zu Prakenstein, Friedrich 27.\,4.\,1867 Meran – 11.\,9.\,1943 ebd.@\textsc{Sölder zu Prakenstein, Friedrich} (27.\,4.\,1867 Meran – 11.\,9.\,1943 ebd.), \emph{Neurologe, Arzt}|pw}{}\ledrightnote{\textcolor{blue}{Friedrich Sölder zu Prakenstein}} Chefarzt; ferner das sog \textcolor{pink}{Maria
                  Theresienschlössel}\oindex{Wien@\textbf{Wien}!XIX., Döbling@\textbf{XIX., Döbling}!Neurologisches Krankenhaus Maria-Theresien-Schlössel@\textbf{Neurologisches Krankenhaus Maria-Theresien-Schlössel}, \emph{Landzunge}|pw}{}\ledrightnote{\textcolor{pink}{Neurologisches Krankenhaus Maria-Theresien-Schlössel}} (\label{K_L04353-2v}\edtext{\textcolor{pink}{Hofstatt}\oindex{Wien@\textbf{Wien}!XIX., Döbling@\textbf{XIX., Döbling}!Hofzeile 18–20@\textbf{Hofzeile 18–20}, \emph{Gebäude}|pw}{}\ledrightnote{\textcolor{pink}{Hofzeile 18–20}}}{\lemma{\textnormal{\emph{Hofstatt}}}\Cendnote{\textnormal{Das Spital lag in der \textcolor{pink}{Hofzeile 18–20}\oindex{Wien@\textbf{Wien}!XIX., Döbling@\textbf{XIX., Döbling}!Hofzeile 18–20@\textbf{Hofzeile 18–20}, \emph{Gebäude}|pwk} und hatte auch einen Trakt, der an der \textcolor{pink}{Pyrkergasse 27}\oindex{Wien@\textbf{Wien}!XIX., Döbling@\textbf{XIX., Döbling}!Pyrkergasse 27@\textbf{Pyrkergasse 27}, \emph{Gebäude}|pwk} lag; bei der \textcolor{pink}{Hofstattgasse}\oindex{Wien@\textbf{Wien}!XVIII., Währing@\textbf{XVIII., Währing}!Hofstattgasse@\textbf{Hofstattgasse}, \emph{Straße}|pwk} handelt es sich um eine
                  Verwechslung.}}}\label{K_L04353-2}, resp. \textcolor{pink}{Pyrkergasse in
                  Döbling}\oindex{Wien@\textbf{Wien}!XIX., Döbling@\textbf{XIX., Döbling}!Pyrkergasse 27@\textbf{Pyrkergasse 27}, \emph{Gebäude}|pw}{}\ledrightnote{\textcolor{pink}{Pyrkergasse 27}}). Prof. \textcolor{blue}{\uline{Redlich}}\pwindex{Redlich, Emil 18.\,1.\,1866 Brünn – 7.\,6.\,1930 Wien@\textsc{Redlich, Emil} (18.\,1.\,1866 Brünn – 7.\,6.\,1930 Wien), \emph{Psychiater, Neurologe, Arzt}|pw}{}\ledrightnote{\textcolor{blue}{Emil Redlich}}; \textcolor{blue}{Julius}\pwindex{Schnitzler, Julius 13.\,7.\,1865 Wien – 29.\,6.\,1939 ebd.@\textsc{Schnitzler, Julius} (13.\,7.\,1865 Wien – 29.\,6.\,1939 ebd.), \emph{Chirurg}|pw}{}\ledrightnote{\textcolor{blue}{Julius Schnitzler}}\textcolor{gray}{w}ill sich an beid\textcolor{gray}{en} Stellen vorerst einmal
               erkundigen; wegen {\pb}Platz und Kosten. Vielleicht rufen Sie
               mich an!\pend
           
\pstart
           Herzlichst Ihr{\\[\baselineskip]}\spacefill\mbox{Arthur}\pend
           \leftskip=0em{}\selectlanguage{ngerman}\endnumbering\briefempfaengerindex{Schwarzkopf, Gustav@\textsc{Schwarzkopf, Gustav}!zzzSchnitzler, Arthur@\emph{von Arthur Schnitzler}!1925-01-041@{4. 1. 1925}|)be}\mylabel{L04353h}
\begin{anhang}
\end{anhang}\newcommand{\dateiname}{L04353}\newcommand{\titel}{Arthur Schnitzler an Gustav Schwarzkopf, 4. 1. 1925}\newcommand{\editorInnen}{Herausgegeben von Jahnke, SelmaMüller, Martin Anton}\input{../tex-inputs/latex-korrekturansicht-abspann}
